\documentclass{arco-article}
\usepackage{arco-acronym}
\usepackage{arco-listings}

\def \note #1 {}
\newcommand{\function}[1]{\texttt{#1}}
\newcommand{\iface}[1]{\textit{#1}}
\newcommand{\id}[1]{\texttt{\textbf{\textscale{.90}{#1}}}}
\newcommand{\fg}[1]{\textit{#1}}
\newcommand{\code}[1]{{\texttt{#1}}}
\newcommand{\seesec}[1]{(véase \S\,\ref{sec:#1})}


\def \mw {middleware}
\def \mws {middlewares}


\usepackage{lettrine}
\newcommand{\drop}[2]{%
  \lettrine[lines=2,findent=2pt,nindent=0pt, loversize=0.25]% lhang=0.33
  {\textcolor[gray]{0.4}{#1}}{#2}%
  }


\title{IDM}

\begin{document}

\note{Protocolos de red incompatibles dificultan la interoperabilidad}
%
\drop{L}{a} gran diversidad en protocolos y tecnologías existentes
dificulta en muchas ocasiones el desarrollo de mecanismos para lograr
interoperabilidad entre aplicaciones que se ejecutan en nodos
cualesquiera, sea cual sea el tipo de red en la que se encuentren.  No
siempre es posible ni deseable utilizar el mismo protocolo de red.

Un buen ejemplo de este problema es el acceso a los nodos de una red
de sensores. Tales dispositivos utilizan tecnologías de red como
Zigbee o Bluetooth, que si bien pueden operar con la pila \ac{TCP/IP},
su adopción suele requerir un esfuerzo considerable y no siempre se
obtienen soluciones óptimas. Otras tecnologías pueden implicar
protocolos especializados o hardware propietario, lo cual en algunos
casos, puede hacer inviable esa posibilidad.

Desde el punto de vista de las \ac{SAN}, éste es un asunto pendiente
~---al menos no totalmente resuelto~--- en lo concerniente al transporte
transparente de mensajes entre una red troncal y las redes de sensores
a las que está conectada. En particular buscamos una solución para los
siguientes dos casos, pero lo suficientemente flexible como para
cubrir muchos otros.

\begin{itemize}
\item Acceso desde el \ac{SI} convencional de cualquier organización a
  las redes de sensores que se desplieguen como soporte a servicios
  cualesquiera.
\item Interacción entre nodos de dos o más redes \ac{SAN} con
  diferente tecnología, incluso si fuera necesario atravesar otras
  intermedias, de cualquier tipo.
\end{itemize}

A fin de cuentas, el objetivo es poder enviar mensajes desde una
red a cualquier otra sin importar que utilicen tecnologías o
protocolos de red incompatibles.


\section{Introducción}

Existen varias alternativas para lograr la interconexión de redes
heterogéneas. En la inmensa mayoría de los protocolos, \mws{} y
sistemas diseñados para la conectividad de redes \ac{SAN} se
considera, casi por descontado, que el sistema cuenta con una única
pasarela de servicio que es utilizada por las redes externas para
acceder a los nodos de la \ac{SAN}. La pasarela aloja normalmente
componentes software cruciales para el acceso a los nodos, lo que
conlleva una gran cantidad de problemas:

\FIXME{que ya se comentaban en la sección~\ref{sec:prev:pasarelas}}

Comparando las diferentes alternativas que se plantean en la citada
sección, parece que la solución más flexible y eficiente es la
utilización de un protocolo de red común, incluyendo la definición de
un esquema de direccionamiento global. Dado que obviamente no se
pretende substituir \ac{IP}, pero sí resolver los problemas que el uso
de la pila \ac{TCP/IP} puede suponer en ciertos ámbitos, el camino
adecuado parece ser un protocolo no solo independiente de la red,
también del transporte. Construir un nuevo protocolo sobre la capa de
enlace implicaría \emph{reinventar la rueda} enfrentándonos a
problemas ya resueltos por otros ya probados y admitidos. La
alternativa más adecuada es usar el protocolo o pila que resulte más
\emph{natural} en cada entorno y construir nuestro protocolo
\emph{homogeneizador} en una capa superior.

La especificación «\ac{CORBA} Message
Routing»~\cite[cap. 17]{corba-messaging} (en adelante \ac{CMR}) apunta
en el camino correcto. El \mw{} proporciona una capa de abstracción lo
suficientemente consistente como para ocultar los detalles de las
tecnologías software y hardware subyacentes, se puede implementar
sobre una gran diversidad de sistemas siempre que se cumplan ciertas
condiciones (en función del \mw{} concreto), proporciona un sistema de
direccionamiento abstracto mediante referencias interoperables a
objetos, etc.

Sin embargo, las motivaciones y decisiones de diseño de \ac{CMR} no
corresponden a un protocolo general de encaminamiento de mensajes. Al
menos los siguientes aspectos resultan inadecuados o contrarios en
relación a los objetivos que perseguimos:

\begin{itemize}
\item Los encaminadores no disponen de tablas de encaminamiento que
  permitan construir rutas extremo a extremo en colaboración con los
  encaminadores vecinos. La ruta queda determinada por las referencias
  indicadas en la propia referencia de objeto destino y/o de
  decisiones administrativas locales del \ac{ORB} en el que ejecuta el
  encaminador intermedio. La ruta que sigue la invocación es una forma
  de lograr tolerancia a fallos ante la imposibilidad de entregar el
  mensaje directamente al destino, o a circunstancias que impliquen la
  desactivación temporal (o definitiva) de parte de los encaminadores
  involucrados.

\item No prevé la posibilidad de implementar encaminamiento dinámico
  que permita establecer rutas alternativas ante cambios en la carga,
  disponibilidad o desempeño de los recursos de la red.

\item Las referencias a objetos, que son el equivalente funcional a
  direcciones universales, pueden tener tamaños arbitrarios. Dado que
  contienen una cantidad no acotada de perfiles y componentes,
  que incluyen a su vez la lista de encaminadores posibles, el tamaño
  de estas referencias puede alcanzar con facilidad miles de bytes.
\end{itemize}

No se puede considerar que éstos sean fallos de diseño en la
especificación de \ac{CMR}. Simplemente ocurre que su finalidad es
diferente, lo cual dificulta su aplicación al problema de la
interconexión genérica de redes.

Este capítulo propone la especificación e implementación de un nuevo
protocolo llamado \ac{IDM}\index{IDM}, y de los mecanismos asociados
para emplear un \ac{MDOO} de propósito general para encapsular y
transportar mensajes como invocaciones a objetos, aprovechando algunos
aspectos del diseño de \ac{CMR}, \ac{IP} y otros protocolos existentes,
pero cuidando otros aspectos que consideramos esenciales para nuestro
propósito:

\begin{itemize}
\item Mantener una sobrecarga lo suficientemente baja como para
  permitir su aplicación a cualquier tipo de red, en especial las que
  tienen un \ac{MTU} pequeño, como es habitual en las \ac{SAN}. Menor
  sobrecarga de cabeceras implica también transmitir menos bytes y por
  tanto consumir menos energía.

\item Incorporación de información tanto a la red como a los propios
  mensajes, que haga posible encaminamiento dirigido por requisitos de
  \ac{QoS}, tanto \fg{DiffServ} (basado en clase) como \fg{IntServ}
  (basado en flujo).

\item El protocolo ha de ser lo suficientemente genérico como para
  implementar algoritmos de encaminamiento dinámico sobre él.

\item Debe tener capacidades de \fg{cross-layer}, que permitan
  eliminar capas o prestaciones innecesarias en función de los
  requisitos demandados por el cliente.

\item Debe definir un esquema de direccionamiento universal,
  jerárquico, con direcciones de tamaño máximo acotado pero lo más
  pequeñas posible. El direccionamiento universal, que potencialmente
  puede incluir nodos de miles de redes, puede implicar direcciones de
  tamaño considerable. Por ello, es necesario llegar a un compromiso que
  satisfaga ambos requisitos contrapuestos.
\end{itemize}

Tanto los elementos finales como los encaminadores de la red se
modelarán e implementarán como objetos distribuidos
convencionales. Además, el programador dispondrá de dos posibles
puntos de vista a la hora de diseñar su aplicación:

\begin{itemize}
\item Un mecanismo transparente para la aplicación cliente, por medio
  de un \ac{API} similar al de los \fg{sockets} \ac{BSD}.
\item Una referencia a un objeto remoto (mediante un \fg{proxy}) sobre
  la que puede invocar métodos conforme a su tipo.
\end{itemize}


\subsection{Inter-redes}

Consideremos el problema como analogía a la red más conocida por
todos. Internet, como en la mayoría de las redes, se fundamenta en
el uso de un protocolo de red único, en este caso \ac{IP}. Un esquema
global de direcciones lógicas permite la interconexión de redes con
distintas tecnologías en la capa de enlace y física.  Por tanto, para
que una red pueda formar parte de una \emph{inter-red} requiere que se
cumplan ciertas condiciones básicas:

\begin{itemize}
\item Los \fg{hosts} deben entender el esquema de direccionamiento lógico.

\item Los hosts deben poder \emph{fabricar} y entramar paquetes sobre
  la tecnología de red subyacente.

\item En el caso de que se trate de un medio de difusión, debe existir
  un mecanismo que permita mapear direcciones lógicas a físicas.

\item Debe existir una pasarela o encaminador capaz de reenviar
  (\emph{forward}) paquetes entre redes diferentes, realizando los
  procesos de desentramado y entramado que correspondan a los
  protocolos de enlace de cada red.
\end{itemize}

En Internet, cuando se precisa atravesar una red que utiliza un
protocolo de red incompatible, la solución suele ser alguna técnica de
\emph{tunneling}, aunque eso impide normalmente implantar un esquema
de encaminamiento homogéneo.  El otro problema evidente es enviar
datos entre dispositivos de esas redes, que tienen protocolos de red y
direcciones lógicas incompatibles.


\subsection{Encaminamiento homogéneo}

Nuestro objetivo es proporcionar un servicio de transporte de mensajes
extremo a extremo, independiente de las tecnologías de las redes
intermedias, aprovechando los protocolos y tecnologías disponibles en
cada dominio a fin de no reimplementar funcionalidades ya
disponibles. En muchos sentidos se trata de cumplir uno de los
requisitos que planteaba el \ac{ISTAG}~\cite{istag01scenarios4ami}
para los entornos inteligentes en cuanto a la necesidad de una
infraestructura de comunicaciones integradora. También es una de las
necesidades identificadas a menudo en los \mws{} para redes de
sensores~\cite{Romer04ProgrammingParadigms}.

\ac{IDM} no delega el direccionamiento a las capas inferiores aunque
sí puede delegar el encaminamiento. En ese aspecto cada red puede
considerarse un \ac{SA}. Por ejemplo, en una red \ac{TCP/IP} se pueden
encapsular los mensajes sobre paquetes \ac{IP} y delegar el
encaminamiento en los protocolos habituales, mientras que en una red
\ac{MANET} sin \ac{IP} se puede implementar un protocolo de
encaminamiento tipo \ac{AODV} cuyos mensajes van encapsulados en
mensajes de nuestro protocolo. La combinación de un \mw{} con la
posibilidad de encaminamiento multicapa (\fg{cross-layer}) extremo a
extremo aporta muchas y valiosas ventajas:

\begin{definitionlist}
\item[Direccionamiento] Realmente global e independiente de la
  tecnología y el protocolo.

\item[Red virtual] Posibilita el despliegue de redes virtuales sobre
  cualquier tipo de red existente o futura.

\item[Calidad de servicio] Proporciona soporte para aplicar políticas
  de calidad de servicio de alto nivel, puesto que se puede manejar
  información mucho más rica en las capas superiores.

\item[Flexible] Y adaptable a las necesidades de la
  aplicación. Utilizar un \mw{} de comunicaciones no implica
  necesariamente una sobrecarga de protocolos ni un aumento
  significativo de la complejidad.  Por ejemplo, es perfectamente
  viable encapsular mensajes sobre tramas Ethernet, si los
  requerimientos de calidad de servicio solicitados lo permiten.

\item[Servicios de la plataforma] Por el hecho de utilizar un \mw{} de
  comunicaciones de propósito general, se puede contar con muchos de
  los servicios comunes que incorpora. En el caso de nuestro prototipo
  actual (basado en \ac{Ice}) podemos destacar: cifrado y autenticación
  transparentes, soporte multi-lenguaje, persistencia automática,
  despliegue, activación implícita, balanceo de carga, replicación
  transparente, etc.

\item[Transparencia de localización] Nuestros encaminadores son
  entidades lógicas y, por tanto, pueden ser migrados o replicados a
  otros nodos de la red de forma transparente.

\item[Encaminamiento funcional] Esta cualidad permite plantear
  esquemas y algoritmos de rutado basados en agrupaciones funcionales
  de nodos y no solo en su localización geográfica.

\item[Eficiencia] Nuestros encaminadores también se pueden implementar
  sobre sistemas empotrados. Es posible utilizar plataformas que van
  desde microcontroladores, como los que incluyen los nodos de una
  \ac{SAN} hasta sistemas sintetizables sobre lógica reconfigurable,
  como \acx{FPGA}. Esto es posible utilizando IceC.

\item[Orientado a Objetos] Si el programador lo desea, es posible ver
  al servidor como un objeto remoto por medio de la declaración de su
  interfaz.  Si se trata de un servicio orientado a flujo o que no
  requiere semántica en la invocación, es posible utilizar un \ac{API}
  alternativo con primitivas \fg{read/write}.
\end{definitionlist}


\subsection{Acerca de los prototipos}

En la sección~\ref{sec:protos:IDM} se describe en detalle el diseño de
algunos de los prototipos iniciales.  Para su implementación se ha
utilizado \ac{Ice}. Sin embargo, las propuestas que se plantean aquí
no son específicas de éste. Se puede utilizar cualquier \mw{}
heterogéneo basado en objetos distribuidos.

A lo largo del texto se harán referencias a algunas de las
características y particularidades de este \mw{} pues consideramos que
puede ayudar a una mejor comprensión de los conceptos que se exponen.


\section{Direccionamiento}

Los esquemas de direccionamiento de red tradicionales utilizan una
dirección lógica que identifica al host, delegando la
identificación de la aplicación destino (el servidor) a los protocolos
de la capa de transporte, llamados \emph{puertos} en el caso de
\ac{TCP} y \ac{UDP}.  \ac{IDM} utiliza un esquema de direcciones
jerárquico, en el que se reserva un sufijo para identificar la
aplicación.  Estas direcciones son realmente identidades de objetos.
Si se desea direccionamiento global, se puede usar cualquier esquema
de identificadores que garantice la unicidad, llamados \acx{GUID}, y
que permita obtener identificadores jerárquicos.  Aunque es posible,
no es estrictamente necesario que la dirección identifique el
host en el que reside el objeto.

El cliente solo necesita conocer la identidad de objeto, que puede
entenderse a todos los efectos como una dirección lógica.  Para llevar
a cabo la transmisión de los datos se precisa de un \fg{endpoint}.
Éste encapsula la implementación del transporte sobre una tecnología o
protocolo concreto: \ac{TCP}, \ac{UDP}, \ac{IP}, Ethernet, \ac{XBow},
Bluetooth, etc.  La elección del endpoint la realiza el
encaminador en base a las restricciones de calidad de servicio
especificadas por el cliente y de la disponibilidad de las interfaces
físicas en el nodo en el que reside.  La transmisión puede no ser
posible si no se satisfacen estas condiciones.

\ac{IDM} proporciona mecanismos para localizar la dirección del
endpoint de un objeto a partir de su identificador sin que sea
necesario conocer su ubicación.  Más concretamente, sirve para
averiguar el endpoint del adaptador en el que está
registrado el objeto destino.


\subsection{Identidad del objeto}

Desde el punto de vista funcional, el identificador de objeto es una
secuencia de bytes globalmente única con dos campos relevantes:

\begin{definitionlist}
\item[Aplicación]

  Identifica el propósito del servicio, es decir, el objetivo común
  de todos los objetos que la comparten.  Ejemplos válidos podrían ser
  una red \ac{P2P}, un conjunto de suscriptores a un servicio de
  \fg{streaming}, etc.  Debe entenderse como un filtro y, por ello, su
  uso no es obligatorio. Este campo puede ser utilizado por los
  encaminadores para dar un tratamiento especial a todos los mensajes
  que corresponden a una aplicación determinada, aplicando técnicas de
  \fg{traffic shaping} o priorización por contenido.

\item[Nombre]

  Identifica el objeto independientemente de la aplicación.  Un objeto
  dado tiene el mismo nombre en todas las aplicaciones en la que esté
  involucrado, aunque también es posible que un mismo objeto tenga
  más de un nombre y, por tanto, varias identidades incluso dentro del
  contexto de la misma aplicación.
\end{definitionlist}

El campo \emph{nombre} debe utilizar necesariamente un esquema
jerárquico, pues esa es la base del funcionamiento de las tablas de
rutas que se emplean.  Una de las características diferenciadoras es
que este tipo de jerarquía puede ser funcional además de geográfica
(ver figura~\ref{fig:functional_routing}). Esto aporta una gran
flexibilidad y permite la migración transparente de nodos y
encaminadores.  Aunque pueden desplegarse encaminadores en cualquier
punto de la red, es imprescindible que existan al menos en las
fronteras entre redes distintas pues el objetivo de \ac{IDM} es
precisamente la interconexión de redes.

\begin{figure}[h]
\begin{center}
\includegraphics[width=0.7\textwidth]{IDM-functional_routing.png}
\caption{\acs{IDM}: Ejemplo de encaminamiento funcional}
\label{fig:functional_routing}
\end{center}
\end{figure}


\subsection{Un posible esquema de direccionamiento}

Para ilustrar el sistema de direccionamiento se propone un esquema
concreto basado en \ac{IPv6}.  Nótese que se trata únicamente de un
ejemplo para ilustrar lo explicado en la sección anterior. El uso de
este tipo de direcciones no implica utilizar ninguna de las
restricciones documentadas en la \ac{RFC} 2460~\cite{rfc2460}. Se
propone como un esquema de direccionamiento general. En principio
\ac{IDM} está diseñado como un protocolo genérico para el transporte
de mensajes en aplicaciones distribuidas, no como un substituto de
\ac{IP}. Normalmente no serán necesarias direcciones tan grandes. En
redes virtuales, de propósito específico o con un número reducido y
acotado de nodos, es posible utilizar direcciones mucho más sencillas,
de sólo unos pocos bytes.

A partir del espacio de direccionamiento de $2^{128}$ definido en
\ac{IPv6}, proponemos un esquema jerárquico multinivel similar al
que se define en la \ac{RFC}2073~\cite{RFC2073}. La
figura~\ref{fig:addressing_scheme} representa ese formato.

\begin{figure}[h]
\begin{center}
\includegraphics[width=0.7\textwidth]{IDM-addr_scheme.png}
\caption{Un posible esquema de direccionamiento para \acs{IDM}}
\label{fig:addressing_scheme}
\end{center}
\end{figure}

El espacio de direcciones queda dividido en los siguientes campos:

\begin{itemize}
\item \textbf{Aplicación} (32 bits). Tal como se indicaba se utiliza
  para identificar una aplicación o red lógica concreta.

\item \textbf{Organización} (32 bits). Identifica la institución
  responsable o propietaria del objeto.

\item \textbf{Sección} (8 bits). Niveles organizativos intermedios.
\item \textbf{Unidad} (8 bits).
\item \textbf{Subred} (16 bits).
\item \textbf{Nodo} (16 bits). Identifica un nodo lógico dentro de la
  red. Un mismo dispositivo puede contener varios.
\item \textbf{Objeto} (16 bits). Identifica un objeto dentro del nodo.
\end{itemize}

Nótese que en ningún momento la utilización de este tipo de
direcciones implica la necesidad de utilizar nada de \ac{IPv6} en la
capa de red, ni tan siquiera una pila de protocolos tipo \ac{TCP/IP}.
Es decir, no existe ninguna posibilidad de ambigüedad o colisión entre
estos identificadores y la dirección lógica de la interfaz de red del
dispositivo en que reside, incluso en el caso en que utilice una
implementación real de \ac{IPv6} en la capa de red.


\section{Endpoints}

Un endpoint es una estructura lógica independiente que encapsula
todos los detalles del transporte de datos para una tecnología o
protocolo subyacente.  Bajo la perspectiva de \ac{IDM}, no hay
diferencias funcionales entre un endpoint que permite abrir y mantener
una conexión \ac{TCP} sobre una \ac{VPN}, de otro que envía tramas
\ac{SLIP} sobre un cable RS-232. Las diferencias están en las
prestaciones del endpoint, pero no en su uso.

El endpoint es responsable de las siguientes tareas:

\begin{itemize}
\item Construye las tramas \ac{IDM}, serializa y encapsula la
  invocación del usuario, calcula sumas de verificación, información
  redundante, autenticación, etc. Todo esto se realiza apoyándose lo
  máximo posible en los servicios proporcionados por el \mw. No tiene
  porqué suponer un trabajo excesivamente complejo.

\item Oculta todos los detalles relacionados con el uso de un
  protocolo o interfaz de comunicaciones.

\item Enumera y comprueba las prestaciones de \ac{QoS} que ofrece cada
  transporte y tecnología concreta.
\end{itemize}

\ac{Ice} incluye de serie únicamente tres tipos de endpoints:
\ac{TCP}, \ac{UDP} y \ac{SSL}.  Nosotros hemos desarrollado algunos
más: Socket UNIX, colas de mensajes y Crossbow.  La
figura~\ref{fig:IDM-layering} ilustra la relación que tiene \ac{IDM}
con el modelo tradicional de capas por medio de los distintos
endpoints. Dicha figura refleja que la elección no está condicionada por
el protocolo de bajo nivel o la capa a la que pertenezca sino por la
funcionalidad que proporcione, es decir, sus prestaciones de \ac{QoS}.

\begin{figure}[h]
\begin{center}
\includegraphics[width=.4\textwidth]{IDM-layering.png}
\caption{Relación de \acs{IDM} con el modelo de capas de red
  tradicional}
\label{fig:IDM-layering}
\end{center}
\end{figure}

Cada endpoint tiene una representación textual que especifica el
protocolo de nivel inferior utilizado, la dirección lógica o física,
el puerto o cualquier otra información necesaria para establecer la
conexión o el envío, y otros datos como \fg{timeouts}, certificados para
cifrado, etc.

En \ac{IDM} existe un tipo de endpoint muy especial que se encarga de
anteponer la cabecera \ac{IDM} a los mensajes del protocolo
inter-\ac{ORB} convencionales.  En el caso de \ac{Ice}, eso simplifica
el uso de \ac{IDM} sobre la distribución software; no hay necesidad de
modificar nada en el código que proporciona el fabricante.

\ac{IDM} podría utilizar el protocolo inter-\ac{ORB} para codificar
sus mensajes (como hace \ac{CMR}) pero ello implica una sobrecarga muy
importante que resulta prohibitiva para \acp{MTU} pequeñas. Por ese
motivo se ha definido un formato de cabecera específico (ver
listado~\ref{code:IDM-format}). Los mensajes de reenvío (\fg{forward})
~---cuyo equivalente en \ac{CMR} sería \function{send\_request()}~--- incluyen
los siguientes campos:

\begin{description}
\item[\id{magic sequence}] (3 bytes) Caracteres `\code{I}',`\code{D}',`\code{M}'.
\item[\id{messageType}] (1 byte) Tipo de mensaje.
\item[\id{messageSize}] (4 bytes) Tamaño total del mensaje incluyendo la
  carga útil.
\item[\id{requestId}] (4 bytes) Identificador de la petición para realizar
  la correspondencia con el mensaje de respuesta si lo hubiere.
\item[\id{routerId}] (2 bytes) Identidad del encaminador que debe
  reenviar este mensaje.
\item[\id{dst}] (16 bytes) La identidad del objeto destino.
\item[\id{src}] (16 bytes) La identidad del objeto que recibirá la respuesta.
\item[\id{hopLimit}] (1 byte) Un contador que el encaminador descuenta en cada
  salto. Evita bucles de encaminamiento de manera similar al campo
  del mismo nombre en {IPv6}.
\item[\id{tos}] (1 byte) Un código que identifica el \ac{ToS}. El
  encaminador aplica un conjunto de restricciones de \ac{QoS} a cada
  tipo de servicio. Su función es similar al campo del mismo nombre en
  \ac{IP}.
\item[\id{flow}] (4 bytes) Identifica una secuencia de mensajes que
  comparten un perfil de \ac{QoS}.
\item[\id{mode}] (1 byte) Características de entrega de la invocación: idempotente,
  \fg{oneway}/\fg{twoway}, respuesta requerida, etc.
\end{description}

\begin{listing}[
  language = Slice,
  caption =  {[Formato de los mensajes \acs{IDM}]Formato de los
    mensajes \acs{IDM}  (notación pseudo-Slice)},
  label =    code:IDM-format]
module IDM {
  struct HeaderData {
    byte magic_I, magic_D, magic_M;
    byte messageType;
    int messageSize;
    int requestId;
    byte routerId;
    IDM::Identity dst;
    IDM::Identity src;
    byte hopLimit;
    ToS tos;
    int flow;
    byte mode;
  };

  struct ForwardData {
    string operation;
    Encapsulation params;
  };

  struct ForwardUserMsg {
    struct HeaderData;
    struct ForwardData;
  };

  struct ForwardRawMsg {
    struct HeaderData;
    Ice::ByteSeq payload;
  };
};

\end{listing}

Nótese que el campo \id{router} no es una identidad \ac{IDM}
completa. Solo se requiere el sufijo ~---la identidad del objeto
dentro del nodo~--- dado que el adaptador que lo aloja debe ser
localizado por otro medio. Por ejemplo, \ac{IP} no requiere
identificar al encargado de procesar los paquetes puesto que
únicamente hay un proceso de encaminamiento por host. Sin embargo, hay
muchos motivos por los que puede resultar ventajoso que varios
encaminadores \ac{IDM} (objetos) convivan al tiempo en un nodo.

A pesar de la mencionada sobrecarga, puede ser interesante definir
este mensaje como un método del encaminador, empleando Slice y
utilizando \ac{IceP} como formato. El caso más claro para decantarse
por esta opción lo tenemos a la hora de prototipar algoritmos de
encaminamiento dinámico o depurar protocolos de aplicación
encapsulados sobre \ac{IDM}. El motivo es que implementar y probar
métodos convencionales del objeto es mucho más sencillo desde el punto
de vista del desarrollo ya que los procesos de (de)serialización son
realizados automáticamente por el \mw.

En cuanto a la carga útil, \ac{IDM} permite transportar mensajes de
dos tipos:

\begin{description}
\item[Raw] Ofrece al usuario un \ac{API} similar a los sockets
  \ac{BSD} para darle la posibilidad de transportar bloques de datos
  sin ningún tipo de semántica ni formato.

\item[User] Si se dispone de la especificación de la interfaz del
  objeto remoto, es posible realizar invocaciones a métodos. De ese
  modo se logra serialización/deserialización de parámetros,
  verificación de tipos, gestión de excepciones, etc.
\end{description}


\section{Gestión de calidad de servicio}

Los requisitos de \ac{QoS} necesarios son especificados en el momento
en que se solicita el acceso al objeto remoto.  Algunos de ellos, como
el cifrado o la autenticación pueden establecerse extremo a extremo,
pero otros como el ancho de banda requerido deben satisfacerse a lo
largo de toda la ruta.  A continuación se listan algunos de los
parámetros que hemos considerado, aunque varios no están disponibles
en los prototipos actuales.

\begin{itemize}
\item Invocación oneway/twoway.
\item Confirmación/Garantía de entrega.
\item Ordenamiento de mensajes.
\item Autenticación.
\item Autorización.
\item Latencia máxima.
\item Ancho de banda mínimo.
\item Encaminamiento fuente.
\end{itemize}

El cliente también debe especificar qué ocurre cuando no es posible
satisfacer alguno de los requisitos. Solicitar una política estricta
puede implicar la imposibilidad de llevar a cabo el transporte.

Se considera que los encaminadores pueden contar con un conjunto de
\ac{ToS} definidos junto con sus requisitos de
\ac{QoS}. Adicionalmente, los encaminadores pueden implementar la
interfaz \iface{IDM::QoS::Admin}, que se utiliza para definir nuevos
perfiles de \ac{QoS}. Esta interfaz (listado~\ref{code:IDM-QoS}) tiene
un método que permite registrar un perfil \ac{QoS} en base a un código
de \ac{ToS} o bien a un identificador de flujo.

\begin{listing}[
  float,
  language = Slice,
  caption =  {Interfaz para administración de perfiles de \acs{QoS} en
              encaminadores \acs{IDM}},
  label =    code:IDM-QoS]
module IDM {
  module QoS {
    enum feature {maxLatency,
                  minBandwidth,
                  reliability};

    dictionary<feature,string> profile;

    interface Admin {
      void setQoSbyToS(profile qos, ToS tos);
      profile getQoSbyToS(ToS tos);
      void setQoSbyFlow(profile qos, int flow);
      profile getQoSbyFlow(int flow);
    };
  };
};
\end{listing}


\section{Entrega de mensajes}

\ac{IDM} se ha concebido como un protocolo de red, de modo que se
plantean los mismos problemas, con la salvedad de que deben resolverse
en una capa de abstracción diferente: sobre el \mw{} en lugar de sobre
la capa de enlace. Un problema básico en un protocolo de red es la
descripción de los mecanismos de entrega de mensajes, tanto a vecinos
de su red como a hosts fuera de ella.


\subsection{Entrega directa}

El método más simple de entrega de mensajes involucra únicamente al
cliente y al objeto o servidor; a esto se denomina normalmente
\emph{entrega directa}.  Para lograrlo es necesario averiguar los
endpoints de que dispone el adaptador al que pertenece el objeto
destino y después comprobar si el host cliente cuenta con algún
endpoint compatible.


\begin{figure}[h]
\begin{center}
\includegraphics[width=0.7\textwidth]{IDM-networks.pdf}
\caption{Una inter-red \acs{IDM} formada por 4 redes \acs{IDM}}
\label{fig:IDM-networks}
\end{center}
\end{figure}

Una \emph{red \ac{IDM}} es un conjunto de objetos directamente
alcanzables por medio de un \emph{tipo} de endpoint.  No tiene
porqué existir ninguna relación entre una red \ac{IDM} y las redes
físicas a las que está conectado el host.  Puede ocurrir incluso que
un adaptador de objetos cuente con varios tipos de endpoints a pesar
de estar conectado a una sola red física. En el escenario más simple
los encaminadores \ac{IDM} se encuentran sólo en los nodos
frontera entre redes con distintas tecnologías o protocolos de red
incompatibles (ver figura~\ref{fig:IDM-networks}). Eso proporciona
conectividad entre todos los objetos y clientes \ac{IDM}, pero hay
otros usos posibles en los que pueden existir muchos más
encaminadores, por ejemplo, en una \ac{MANET} todos los nodos son
encaminadores \ac{IDM}.

Para conocer los endpoints del adaptador al que pertenece un objeto se
ha diseñado el protocolo \ac{ALP}. Es un protocolo general de
descubrimiento de vecinos, pero asume que los actores son objetos
accesibles mediante \ac{IDM}.  El listado~\ref{code:ALP} muestra la
especificación de la interfaz que deben implementar los objetos que
soportar este protocolo.

\note{idm/slice/ALP.ice}
%
\begin{listing}[
  float,
  language = Slice,
  caption =  {[Interfaz \acs{ALP} para registro de objetos y localización de
             adaptadores] Interfaz \acs{ALP} para registro de objetos y
             localización de adaptadores (especificación Slice)},
  label =    code:ALP]
module IDM {
  module ALP {
    exception ObjectNotFoundException {};

    interface Locator {
      void add(Object* prx);
      void remove(IDM::Identity oid);

      Object* resolve(IDM::Identity oid)
        throws ObjectNotFoundException;
    };

    interface Lookup {
      idempotent void request(IDM::Identity oid,
                              Lookup* callback);
      idempotent void reply(Object* prx);
    };
  };
};
\end{listing}


% Locators y ALP
Cada nodo de la red \ac{IDM} debe disponer de un \emph{demonio} local,
llamado \emph{Servicio \ac{ALP}}, que implementa la interfaz
\iface{ALP::Locator}. Mediante la invocación del método
\function{lookup} y conocida la identidad destino, el nodo cliente
obtiene un proxy, que permite realizar invocaciones directamente
sobre el objeto remoto. Nótese que sólo es posible localizar endpoints
de adaptadores vecinos, es decir, que se correspondan a nodos de la
misma \emph{red \ac{IDM}}.

Existen diversas formas de implementar el servicio \ac{ALP}. Para
nuestros prototipos iniciales hemos considerado dos. La elección se
hace por configuración y es transparente tanto para los clientes como
para los objetos.

\begin{itemize}
\item Un servicio distribuido que permite a cualquier host interrogar
  a su segmento de red sobre la existencia de un objeto concreto. Este
  servicio utiliza los mecanismos \fg{multicast} o \fg{broadcast}
  específicos de los endpoints disponibles. El servicio en cada
  nodo debe atender los mensajes \function{request()} y los debe
  contestar mediante mensajes \function{reply()}. Estos mensajes
  pertenecen a la interfaz \emph{IDM::ALP::Lookup} y en la práctica
  dicha interfaz se implementa como una faceta del \emph{Servicio ALP}
  (ver figura~\ref{fig:IDM-blocks}).

\item Utilizar las características de \emph{indirección} que incorpore
  el \mw. En \ac{Ice} la indirección de objetos se consigue por medio
  de un servicio llamado \emph{Ice Locator}, que aunque centralizado
  puede estar replicado. \ac{IDM} dispone de un adaptador para
  utilizar \iface{IceGrid::Locator} con una sobrecarga mínima.  Eso
  resulta muy útil en segmentos grandes de la red virtual y cuando los
  hosts dispongan de suficientes recursos de cómputo como para
  ejecutar un nodo IceGrid \seesec{prev:icegrid}. Obviamente ofrece
  alguna desventaja respecto al anterior: todos los hosts deben
  conocer desde su arranque la referencia remota de al menos una de
  las réplicas del \emph{Locator} y contar con un endpoint
  compatible con ella.
\end{itemize}

\begin{figure}[h]
\begin{center}
\includegraphics[width=.7\textwidth]{IDM-blocks.png}
\caption{Diagrama de bloques de \acs{IDM}}
\label{fig:IDM-blocks}
\end{center}
\end{figure}


En ambos casos, los hosts que deseen exportar objetos deben
registrarlos en el servicio \ac{ALP} para que puedan ser encontrados
por otros nodos, independientemente si están implementados mediante
\emph{IceGrid Locator} o por el propio host.  Ambos mecanismos pueden
coexistir en la misma red.

El servicio \ac{ALP} cuenta con una caché que almacena temporalmente
los endpoints resueltos, sin importar el modo en que se
resolvieron.  El propósito de la caché es reducir el impacto en el
desempeño que implican los mecanismos de resolución. Cuando se utiliza
un servicio \ac{ALP} centralizado, el \fg{runtime} \ac{IDM} incorpora
una caché local para evitar tener que repetir consultas al servicio
remoto.


\subsection{Entrega indirecta}

Se utiliza entrega indirecta para aquellos objetos que no son
alcanzables a través de ninguno de los endpoints del cliente.  En
esos casos es necesaria la mediación de al menos un encaminador.
Nótese que puede ocurrir un caso especial.  Aún existiendo
accesibilidad directa por medio de un endpoint, es posible que
las restricciones de \ac{QoS} impidan su uso y obliguen a utilizar uno
o varios encaminadores para llegar al destino.


\section{Encaminamiento}

Todos los nodos que forman parte de una inter-red \ac{IDM} requieren
de una tabla de rutas que determina si la entrega de mensajes puede
realizarse de forma directa o se necesita un encaminador. Éste también
es un objeto distribuido convencional, similar a los nodos servidores; la
principal diferencia es que puede aceptar y redirigir mensajes que no
van dirigidos explícitamente a él.

El listado~\ref{code:IDM-rute-table} muestra la tabla de rutas de un
encaminador. Se han utilizado identificadores textuales simplificados
para mejorar la legibilidad.  Esta tabla de rutas muestra lo
siguiente:

\begin{itemize}
\item Los objetos con prefijo {\small\texttt{NetA.Subnet1}} o
  {\small\texttt{NetA.Subnet2}} son alcanzables directamente (entrega
  directa).

\item Los objetos con prefijo {\small\texttt{NetB}} son alcanzables a través
  del encaminador {\small\texttt{NetA.Subnet1.R2}}.

\item Para el resto de objetos se debe utilizar el encaminador
  {\small\texttt{NetA.Subnet2.R3}} (ruta por defecto).
\end{itemize}

\begin{listing}[
  float,
  caption = Tabla de rutas IDM simplificada,
  label =   code:IDM-rute-table]
Router.Row.NetA.Subnet1.* = tcp, *
Router.Row.NetA.Subnet2.* = tcp, *
Router.Row.NetB. *        = udp tcp, NetA.Subnet1.R2
Router.Row.*              = tcp, NetA.Subnet2.R1
\end{listing}

Cada fila de la tabla de rutas incluye una lista de tipos de
endpoints que se pueden utilizar para contactar con objetos de
esa red.  Eso se debe a que cada red \ac{IDM} está constituida por una
o más redes físicas compatibles que utilizan un mismo protocolo (o
conjunto de ellos).  Para utilizar los mecanismos \ac{ALP} de una red
\ac{IDM} es necesario conocer al menos un tipo de endpoint compatible
con esa red.

La figura~\ref{fig:IPvsIDM} muestra las similitudes funcionales
(salvando las distancias) entre los procesos de encapsulación en un
encaminador \ac{IP} y uno \ac{IDM}.

\begin{figure}[h]
\begin{center}
\includegraphics[width=.6\textwidth]{IPvsIDM.png}
\caption{Comparación entre un encaminador \acs{IP} y uno \acs{IDM}}
\label{fig:IPvsIDM}
\end{center}
\end{figure}



\subsection{Encaminamiento dinámico}

Como los encaminadores son objetos distribuidos, resulta relativamente
sencillo implementar algoritmos de encaminamiento dinámico de
cualquier tipo.  Cada mensaje del protocolo de encaminamiento elegido
se define como un método del objeto en la interfaz específica.  Para
soportar un nuevo protocolo de encaminamiento es suficiente que
implemente la interfaz correspondiente como una nueva faceta.

Como ejemplo, el listado~\ref{code:RIP} muestra la interfaz para el
protocolo \ac{RIP}2~\cite{rfc2453}. En ambos métodos, el parámetro
\id{oids} es una lista de identificadores de objetos, que puede estar
vacía en el mensaje \function{request()}. El parámetro \id{src} es la
identidad del encaminador que realiza la invocación.

Este planteamiento permite implementar varios protocolos de
encaminamiento dinámico en el mismo nodo, aunque de operar a la vez,
debe definirse un criterio de prioridad para evitar conflictos en la
actualización de la tabla de rutas.

\note{idm/slice/Routing.ice}
%
\begin{listing}[
  float,
  language = Slice,
  caption =  {Interfaz Slice para el protocolo de encaminamiento \acs{RIP}},
  label =    code:RIP]
module Routing {
   sequence <Ice::Identity> IdentitySeq;

   interface RIP {
      void request(IdentitySeq oids, Ice::Identity src);
      void response(IdentitySeq oids, Ice::Identity src);
   };
};
\end{listing}


\subsection{Gestión de la red}

Dado que los encaminadores y los servidores son objetos distribuidos
convencionales, es posible utilizar los servicios estándar para
despliegue que proporciona el \mw{} (IceGrid en \ac{Ice}).
Con ello se consiguen muchas funcionalidades, algunas típicas de una
plataforma de gestión de red, pero con sobrecarga mínima y totalmente
integradas en el funcionamiento de la propia red:

\begin{definitionlist}
\item[Planificación]

  Es posible diseñar la topología, analizar las necesidades y
  requerimientos de red y los usuarios, cuantificar los costes de
  mantenimiento y operación incluso antes de disponer físicamente de
  los hosts y el equipamiento de comunicaciones.

\item[Despliegue]

  Permite realizar distribución de ficheros de configuración,
  bibliotecas y binarios a cada encaminador y, todo ello, a la medida
  de las prestaciones y tipo de arquitectura de cada nodo.

\item[Monitorización/contabilidad]

  Se puede evaluar el rendimiento del sistema, consecución de las
  prestaciones, ancho de banda disponible/consumido, latencia y muchas
  otras variables de un modo sencillo en cada encaminador y servidor
  de la red.

\item[Gestión de la configuración]

  Cuando se disponga de nuevas versiones del software, éste puede ser
  desplegado a los encaminadores, incluyendo implementación o mejoras
  de los protocolos de encaminamiento. También es sencillo añadir
  interfaces para administración de los equipos, definición de
  prioridades, políticas de seguridad, etc.
\end{definitionlist}


\section{Integración con el \mw}

Una de nuestras prioridades es la integración del mecanismo de
encaminamiento con el \mw. Para el programador, las diferencias
entre usar un objeto \ac{IDM} y un objeto distribuido convencional
deberían ser mínimas. Por ese motivo, el reenvío de mensajes se
utiliza únicamente entre encaminadores y no afecta a la interacción de
estos con los clientes y los objetos.

Cuando un cliente desea hacer un invocación a un objeto \ac{IDM}
remoto utiliza una clase especial para conseguir un proxy a su
encaminador \emph{local} (situado en el mismo nodo que el cliente). El
cliente realiza la invocación utilizando este proxy como si se tratase
realmente del objeto remoto.

A continuación se explica en detalle el proceso completo para realizar
una invocación a un objeto \ac{IDM} remoto afectando al menos a un
encaminador intermedio, es decir, potencialmente cliente y objeto
destino podrían encontrarse en redes tecnológicamente distintas. La
figura~\ref{fig:IDM-invocation} ilustra todo ese proceso.

\begin{figure}[h]
\begin{center}
\includegraphics[width=\textwidth]{IDM-invocation.png}
\caption{\acs{IDM}: Invocación de un objeto \acs{IDM} remoto}
\label{fig:IDM-invocation}
\end{center}
\end{figure}

El encaminador local puede recibir invocaciones destinadas a cualquier
objeto. Si la identidad destino es la suya propia, él procesa el
mensaje; en caso contrario, encamina el mensaje hacia su destino.
Cuando un encaminador recibe una invocación de un cliente local
realiza los siguientes pasos:

\begin{enumerate}
\item Si el mensaje requiere respuesta, crea un objeto \fg{callback}
  alojado en su espacio de memoria para poder recibir la respuesta. La
  referencia a este objeto es la que aparecerá como campo \id{src} de
  la estructura \id{HeaderData}.
\item Construye un mensaje según el formato de la estructura
  \id{ForwardRawMsg} o \id{ForwardUserMsg} (en adelante \fg{forward},
  por simplicidad) según el tipo de invocación elegida por el
  cliente.
\item Comprueba su tabla de encaminamiento y determina la identidad
  del siguiente encaminador y los endpoints de que dispone.
\item Filtra dichos endpoints según la información de \ac{QoS} de cada
  uno de ellos en relación al \ac{ToS} al que corresponde la
  invocación procedente del cliente. Si no encuentra ninguno
  compatible retorna una excepción al cliente y aborta.
\item Utiliza el servicio \ac{ALP} para determinar la información de
  direccionamiento del protocolo inferior que corresponda.
\item Utiliza dicho endpoint para enviar el mensaje \fg{forward} al
  siguiente encaminador.
\end{enumerate}

Encaminador \emph{intermedio} es aquel que recibe mensajes
\fg{forward} procedentes de cualquier otro encaminador. En este caso
el procedimiento de trabajo incluye únicamente los puntos 3 a 6 de la
lista anterior.

El encaminador \emph{final}  debe encargarse de entregar invocaciones
directamente a los objetos destino. Procede del siguiente
modo:

\begin{enumerate}
\item Recibe un mensaje \fg{forward} y comprueba su tabla de rutas
  para determinar si el destino es un objeto de la forma habitual.
\item Mediante \ac{ALP} obtiene los endpoints del objeto destino. Si
  ninguno de ellos fuera compatible con los requisitos de \ac{QoS}, en
  este momento generaría una excepción dirigida al objeto \id{src}.
\item A partir del contenido del mensaje \fg{forward} crea un mensaje
  de invocación convencional (en el caso de \ac{Ice} será un
  \id{HeaderData} y un \id{RequestData} del protocolo \ac{IceP}).
\item Ejecuta dicha invocación directamente sobre el cliente. Puede
  ser síncrono o no dependiendo del tipo de mensaje y el endpoint
  utilizado.
\item Si la invocación es síncrona y hay respuesta, el encaminador
  procede con los mismos pasos que se han descrito para el encaminador
  local, pero utilizando la referencia \id{src} contenida en el
  mensaje \fg{forward}.
\end{enumerate}

Cuando la respuesta llega al objeto temporal que creó el encaminador
local, éste crea una respuesta según el formato del protocolo
inter-\ac{ORB} y lo entrega al cliente por la misma conexión que se
estableció para realizar la invocación. El cliente ha estado bloqueado
durante todo el proceso a menos que utilice los mecanismos \ac{AMI}
proporcionados por el \mw.

Los calificativos \emph{local}, \emph{final} e \emph{intermedio} no
determinan tipos de encaminadores sino roles. Todos los encaminadores
pueden funcionar con cualquiera de esos roles. La distinción
presentada aquí se debe a motivos de legibilidad y
simplificación de las fases del procedimiento.

Nótese que a diferencia de \ac{CMR}, se asume que todo nodo que aloja
clientes tiene un encaminador local, que es el único que recibe
invocaciones directamente de los clientes. Es decir, un cliente no
puede efectuar invocaciones convencionales (no-\ac{IDM}) a
encaminadores fuera de su nodo. Aunque esto implica un salto más,
mejora la escalabilidad del sistema dado que los encaminadores
tienen que almacenar objetos callback para sus clientes. De ese
modo, es mucho más sencillo dimensionar las necesidades de memoria y
procesamiento del nodo, y los encaminadores sin clientes ni objetos
pueden ser mucho más sencillos pues están por diseño exentos de esa
tarea.


\subsection{API}

La intención es que \ac{IDM} esté disponible como una biblioteca
para servicios avanzados de comunicaciones extremo a extremo de
forma transparente.  Por ese motivo, se ha diseñado e implementado
como un conjunto de clases que permite desarrollar objetos y
clientes de una manera muy sencilla.

Desde el punto de vista del programador de aplicaciones, para la
implementación de servidores sólo es necesario tener en cuenta una
pequeña diferencia respecto a la implementación de un servidor
\ac{Ice} convencional. Se debe utilizar un \emph{Adapter} especial,
llamado \iface{IDM::Adapter}. Es un \emph{decorador}\footnote{en
  referencia al patrón de diseño
  \fg{decorator}~\cite{design_patterns}} que extiende la funcionalidad
del \emph{Adapter} \ac{Ice}.  Al crearse este adaptador, se registra
automáticamente en el \emph{Servicio ALP} disponible (tal como se haya
especificado por configuración). El listado~\ref{code:IDM-server}
muestra la implementación en lenguaje C++ de un objeto mínimo
accesible mediante \ac{IDM}.

\lstinputlisting[
  language = C++,
  caption =  {Implementación de un servidor \acs{IDM} mínimo},
  label =    code:IDM-server]
{listings/idm-server.cpp}

El cliente \ac{IDM} también requiere algunas modificaciones respecto a
la implementación de un cliente \ac{Ice} convencional.  Es necesario
crear una instancia de la clase \iface{IDM::Socket}.  Éste dispone de
un método \function{plug()} que acepta un identificador de objeto y
devuelve un proxy que permite enviar mensajes a ese objeto.
Internamente, esa clase consulta la tabla de encaminamiento del nodo y
utiliza el servicio \ac{ALP} para conseguir la referencia ya sea para
entrega directa o indirecta (primer salto).  El
listado~\ref{code:IDM-client} muestra la implementación de un cliente
para acceder al servidor del listado~\ref{code:IDM-server}.

\lstinputlisting[
  float,
  language = C++,
  caption =  {Implementación de un cliente \acs{IDM} mínimo},
  label =    code:IDM-client]
{listings/idm-client.cpp}


\section{Conclusiones}

Aunque a primera vista, el diseño e implementación de un protocolo de
red por encima incluso de la capa de transporte pueda parecer un
concepto contradictorio y un esfuerzo baldío, queda demostrado que las
prestaciones que ofrece \ac{IDM} pueden resultar muy valiosas en
muchos escenarios diferentes.  Algunas de las más importantes son:

\begin{itemize}
\item Direccionamiento lógico universal independiente de la tecnología.
\item Transporte transparente extremo a extremo.
\item Mejoras en la eficiencia al poder seleccionar transportes
  específicos con menor granularidad, que decide el diseñador de la
  aplicación.
\item Prestaciones adicionales transparentes y desacopladas del
  transporte como cifrado, contabilidad, autenticación, autorización,
  etc.
\item Encaminamiento funcional.
\item Migración de nodos y encaminadores.
\end{itemize}

Aprovechando las posibilidades que ofrece \ac{IDM}, estamos evaluando su
aplicación a diversos campos:

\begin{itemize}
\item Redes \acx{MANET}. Cada nodo de la red ad hoc puede incluir un
  encaminador \ac{IDM}, homogeneizando el acceso desde la red
  troncal. Este enfoque permite incluso plantear redes ad hoc
  híbridas como, por ejemplo, una red \ac{IP} y una red de sensores.

\item Implementación del encaminador \ac{IDM} en sistemas empotrados
  basados en microcontrolador y/o en \ac{FPGA} para el desarrollo de
  encaminadores eficientes para redes \ac{SAN} y grids
  heterogéneos~\cite{Rincon09SupportingOS}.

\item Métricas del desempeño de la red y del servicio prestado.

\item Red virtual para provisión de \ac{QoS}~\cite{Urzaiz09Anovel}.
\end{itemize}

En lo referente al desarrollo de nuevos endpoints, cabe destacar
que se encuentran en fase de desarrollo los que corresponden a
protocolos tan variados como Zigbee, BlueTooth y Ethernet.

Un concepto que creemos merece la pena abordar es lo que llamamos
\emph{decoradores}~\footnote{También en referencia al patrón de
  diseño.} de endpoints.  Pueden añadir características opcionales a
cualquier endpoint, como cifrado, garantía de entrega, contabilidad,
etc.  Estos decoradores facilitan la composición funcional de
endpoints dando lugar a combinaciones muy interesantes.  Por
ejemplo, si se requiere transmitir audio en una red \ac{LAN} Ethernet
salvaguardando la privacidad, pero sin necesidad de garantizar la
entrega, se podría utilizar un endpoint que realiza transporte
directamente sobre tramas Ethernet y que está \emph{decorado} con un
endpoint que cifra los mensajes.  De ese modo, se puede conseguir el
transporte deseado sin la sobrecarga que implicarían las cabeceras
\ac{IP} y \ac{UDP}.



\end{document}

% Local Variables:
%   coding: utf-8
%   fill-column: 90
%   mode: flyspell
%   ispell-local-dictionary: "castellano8"
%   mode: latex
% End:

%  LocalWords:  interoperabilidad middleware
